%--------------------------------------------------------------------------------
%Plantilla de latex - formato para tesis, IngenierÍa Civil UNSCH.
%Copyright (C) 2025 Cristhoper Jairo Flores Peña.
%--------------------------------------------------------------------------------
\documentclass[10pt,a4paper,oneside]{book}
\usepackage{EstiloTesis}
\usepackage{bookmark}%bookmark: Mejora la navegación en el PDF generado con marcadores interactivos.
\usepackage{subfigure}%subfigure: Permite incluir subfiguras dentro de una figura.
\usepackage{float} %float: Mejora el control de posicionamiento de figuras y tablas.
\usepackage{multirow}%Permite crear celdas que abarcan múltiples filas en tablas.
\usepackage{microtype}%Optimiza el espaciado y la tipografía para mejorar la legibilidad.
\usepackage{fontenc}%Define la codificación de fuentes (generalmente usado con T1 para idiomas con acentos).
\usepackage{amssymb}%Proporciona símbolos matemáticos adicionales
\usepackage{pgfplots}
\pgfplotsset{compat=1.15}


\definecolor{rojo}{RGB}{255,0,0}
\definecolor{amarillo}{RGB}{255,255,0}
\definecolor{blanco}{RGB}{255,255,255}



%-----------------------------------------------------------
\addbibresource{Capitulos/Referencias.bib}
%--------------------------------------------------------------------------------

%  MACROS
%--------------------------------------------------------------------------------
\universidad{UNIVERSIDAD NACIONAL DE SAN CRISTÓBAL DE HUAMANGA}
\facultad{FACULTAD DE INGENIERÍA DE MINAS, GEOLOGÍA Y CIVIL}
\escuela{ESCUELA PROFESIONAL DE INGENIERÍA CIVIL}
\tema{\textbf{"Desarrollo de una arquitectura para la evaluación de baches en pavimentos flexibles mediante el análisis de representaciones de profundidad con aprendizaje profundo, Ayacucho 2025"}}

\autor{\textbf{Bach. Cristhoper Jairo Flores Peña}}
\asesor{\textbf{Mg. Ing. Alex Sander Ircañaupa Huamani}}

\AtEveryCite{\defcounter{maxnames}{1}}
\begin{document}


\selectlanguage{spanish}
%\thess % INCLUYE: INDICE GENERAL, PARTE, APENDICE FIGURA, CAPITULO, ETC.)
\frontmatter    % INCLUYE: TITULO, PROLOGO
\pagestyle{plain} %establece el estilo de página para todo el documento
% CUERPO PRELIMINAR(I)
\include{Capitulos/0_1Portada}
%\include{Capitulos/0_2Derechos}
%\include{Capitulos/0_5Resumen}
%\include{Capitulos/0_31Presentacion}
%\include{Capitulos/0_3Dedicatoria}
%\include{Capitulos/0_4Agradecimiento}
\tableofcontents    % ÍNDICE DEL CONTENIDO
%\newpage
%\listoffigures      % ÍNDICE DE FIGURAS
%\newpage
%\listoftables 		% ÍNDICE DE TABLAS
%\newpage                     		
%\pagestyle{Glosario}
%\include{Capitulos/0_7Glosario}
%\pagestyle{acronimos}
%\include{Capitulos/0_6Nomenclatura}

%--------------------------------------------------
\pagestyle{Normal}
\mainmatter     % INCLUYE LOS CAPÍTULOS
\include{Capitulos/1_CapituloI} 		%I	.-PLANTEAMIENTO DEL ESTUDIO
\include{Capitulos/2_CapituloII}		%II	.-MARCO TEÓRICO
\include{Capitulos/3_CapituloIII}		%III.-MATERIALES Y METODOS
\include{Capitulos/4_CapituloIV}		%IV	.-RESULTADOS
%\include{Capitulos/5_CapituloV}         %V  .-APRECIACIONES
%\include{Capitulos/5_Conclusiones}		%CONCLUSIONES
%\pagestyle{Conclusiones}
%\include{Capitulos/5_Conclusiones}		%CONCLUSIONES

%%% BIBLIOGRAFÍA

\newpage
\pagestyle{RB}
\medskip
\nocite{*}
\printbibliography
%\bibliographystyle{apacite}
%\bibliography{Capitulos/Referencias}
\printbibheading[heading=bibintoc,title={REFERENCIAS BIBLIOGRÁFICAS}]
\printbibliography[title={Referencias Bibliográficas}]
%\printbibheading[heading=sybibintoc]
%\printbibliography[heading=sysubbibintoc,keyword={CON},title={Concreto}]
%\printbibliography[heading=subbibintoc,keyword={OPTI},title={Optimizar Concreto}]
%\printbibliography[heading=subbibintoc,keyword={TEC},title={Tecnología}]
%\printbibliography[heading=sysubbibintoc,keyword={MAN},title={Libros y Reglamentos}]
%\printbibliography[heading=subbibintoc,keyword={MET},title={Método de Investigación}]

%--> Bibliografia 2 ( bibtex % )
%\nocite{*}
%\bibliographystyle{unsrt}
%\bibliography{Capitulos/Referencias}
%\pagestyle{Normal}-zzz
%\mainmatter-zzz
%\appendix-zzz
%\renewcommand{\thefigure}{\Alph{chapter}.\arabic{figure}}
%\renewcommand{\thetable}{\Alph{chapter}.\arabic{table}}
%\include{Capitulos/Apendices}
%\include{Capitulos/Apendices2}

\backmatter     % INCLUYE LOS CAPÍTULOS SIN ENUMERARLOS
\end{document} 